\subsection{30}
30 er et terningespil der spilles med 6 terninger, hvor det gælder om at slå over $\geq 30$ eller $\leq 10$.

Til spillet skal der bruges: 1 raflebæger og 6 terninger. Man starter med at slå med alle 6 terninger, og tager minimum én terning fra\footnote{Man SKAL tage én terning hver gang man laver et slag.}. Man slår da igen med de 5 (eller færre) terninger og igen tager minimum én terning fra. Dette fortsætter indtil man ikke kan tage flere terninger fra, hvorefter man summer alle terningernes øjne sammen.

Afhængigt af terningernes sum er man nu i én af følgende situationer:
\begin{itemize}
    \item Hvis man har slået $<30$ så skal man drikke tåre op til $30$. F.eks. hvis der er en person der har slået $4 + 5 + 6 + 4 + 2 + 6 = 27$ så skal de drikke 3 tåre.
    \item Hvis man har slået præcis $30$ så er der fællesskål og alle drikker.
    \item Hvis man har slået $>30$ så skal man slå efter det tal man har over $30$.
\end{itemize}

Følgende er et eksempel på den sidste situation: Så hvis man har slået $32$, slår man efter $2$'er, eller hvis man har slået $34$ skal man slå efter $4$'er. Igen slår man med alle $6$ terninger og, hvis man slår det tal man er efter eksempelvis $2$'er så tager man dem fra og slår igen med de resterende terninger indtil man ikke slår flere $2$'er\footnote{Eller hvilket tal man nu er efter.}. Man summer igen terningernes øjne, og denne sum er det antal tåre man må give ud til medspillere. Tårene skal gives ud som heltals multiplum af det man har slået, for eksempel: Hvis en person har slået tre $4$'er, kan vedkomne uddele
\begin{itemize}
    \item 4 tåre til tre personer,
    \item 8 tåre til en person og 4 til en ande,
    \item eller 12 til en enkelt person.
\end{itemize}

Hvis man skal slå efter et tal over $30$, og man i det første slag med alle $6$ terninger IKKE slår det tal man er efter, så slå man mod sig selv. Her er reglerne det samme som før, dog med den ændring at man selv skal drikke alle de tåre man slår. Så i det ovenstående tilfælde ville man skulle drikke $12$ tåre selv.

Dette "frem og tilbage" mellem en selv og ens modstandere fortsætter indtil første gang der bliver slået det tal man er efter.

Derefter er det den næste persons tur og man spiller indtil alle er gået kolde eller man ikke gider mere.


